\section{Évaluation de l'option sur OAT}
\label{sec:evaluation}

\subsection{Le contrat}

L'option européenne d'échéance $T_0$ et de strike $K$, exprimé sur le prix clean, paie
$\big(\bar B^{clean}(T_0) - K\big)^+$ pour un call et $\big(K - \bar B^{clean}(T_0)\big)^+$ pour un put. Le
coupon couru $CC(T_0)$ à l'échéance est connu dès aujourd'hui, si bien que
$\bar B^{clean}(T_0) - K = \bar B(T_0) - X$ avec le strike dirty $X = K + CC(T_0)$ : le modèle travaille sur le
prix dirty, somme actualisée des flux. L'option est collatéralisée et actualisée au taux €STR. Sous la mesure
$T_0$-forward €STR $\Q^{T_0}$, de numéraire $P(\cdot,T_0)$ \citep[CM3]{Hillairet2026}, son prix est
\begin{equation}
C = P(0,T_0)\,\E^{T_0}\big[(\bar B(T_0) - X)^+\big].
\label{eq:prix_call}
\end{equation}
Les strikes sont exprimés en pourcentage du forward clean défini ci-dessous, et les prix en pourcentage du nominal.

\subsection{Le forward repo}

Acheter l'OAT au prix dirty de marché $\bar B^M(0)$, la financer en repo jusqu'en $T_0$ et rembourser le
financement avec les coupons encaissés d'ici là permet de livrer le titre en $T_0$. Par absence d'arbitrage, le
prix forward dirty est
\begin{equation}
F = \frac{\bar B^M(0) - \sum_{0 < T_i \le T_0} K_i\,P^{repo}(0,T_i)}{P^{repo}(0,T_0)},\qquad
P^{repo}(0,T) = P^M(0,T)\,e^{-s^{repo}(T)T},
\label{eq:forward}
\end{equation}
où $s^{repo}(T)$ est le spread du taux repo zéro continu contre l'€STR. Un contrat forward collatéralisé en €STR
a pour prix d'équilibre $\E^{T_0}[\bar B(T_0)]$. Comme l'écart entre repo et €STR est déterministe, le portage se
réplique, et $\E^{T_0}[\bar B(T_0)] = F$. C'est la relation forward d'un actif stockable
\citep[cours~2, slides~16-18]{MPR2025} : vu de l'€STR, le spread repo est l'opposé d'un rendement d'opportunité
(\emph{convenience yield}), qui serait positif pour un titre «~spécial~» en repo. Le forward clean, auquel on
rapporte les strikes, est $F^{clean} = F - CC(T_0)$, et la parité call-put s'écrit
$C - \text{Put} = P(0,T_0)\,(F^{clean} - K)$.

Le forward n'utilise le spread repo qu'aux dates antérieures à l'échéance : la courbe repo doit couvrir $T_0$.
Elle est définie jusqu'à dix ans, la plus longue échéance étudiée, et n'est jamais prolongée, de sorte qu'aucun
prix ne repose sur une extrapolation du financement.

\subsection{Recentrer le modèle sur le forward repo}

Laissé à lui-même, le modèle fait croître le prix de l'OAT au taux $\bar r$ en moyenne. Son espérance forward
s'écrit
\begin{equation}
\begin{split}
M &= \sum_{T_i>T_0} m_i,\\
m_i &= K_i\,\E^{T_0}\big[\bar P(T_0,T_i)\big]
= K_i A_i\exp\Big(-H_x^i\mu_x - H_y^i\mu_y + \tfrac12\Var\big(H_x^i x(T_0) + H_y^i y(T_0)\big)\Big),
\end{split}
\end{equation}
où $A_i$ est la partie déterministe de $\bar P(T_0,T_i)$ dans~\eqref{eq:zc_oat}, $H^i_\cdot = H_{a_\cdot}(T_0,T_i)$,
et $(\mu_x,\mu_y)$ les moyennes de $(x(T_0),y(T_0))$ sous $\Q^{T_0}$. Comme dans
\citet[lemme~4.2.2]{BrigoMercurio2006}, le changement de numéraire ajoute à la dérive des facteurs leur covariance
avec le numéraire. Celui-ci est ici le zéro-coupon €STR, qui ne dépend que de $x$ ; en intégrant,
\begin{equation}
\mu_x = -\frac{\sigma_x^2}{2a_x^2}\big(1-e^{-a_xT_0}\big)^2,\qquad
\mu_y = -\frac{\rho\sigma_x\sigma_y}{a_x}\Big[\frac{1-e^{-a_yT_0}}{a_y} - \frac{1-e^{-(a_x+a_y)T_0}}{a_x+a_y}\Big],
\end{equation}
soit $-\Cov\big(x(T_0),\int_0^{T_0}x\big)$ et $-\Cov\big(y(T_0),\int_0^{T_0}x\big)$ ; variances et covariance
sont inchangées.

Le porteur finance l'OAT au taux repo, et non à son propre taux : $M$ ne coïncide pas avec $F$, et l'écart croît
avec $T_0$ comme le portage accumulé. On recale donc le modèle sur le marché, comme $\bar\alpha$ recale la
courbe, en décalant la moyenne du facteur spread en $T_0$ d'une constante $\delta$ telle que
\begin{equation}
\sum_{T_i>T_0} m_i(\delta) = F,\qquad m_i(\delta) = m_i\,e^{-H_y^i\,\delta}.
\label{eq:recentrage}
\end{equation}
$\delta$ s'obtient par une résolution numérique à une dimension d'une expression fermée.

À un terme de convexité près, $\delta$ est nul quand le spread repo égale le spread de l'OAT : le portage est
alors celui de la courbe OAT, et $M = F$. Quand le spread repo est plus faible, le porteur financé en repo
encaisse un portage positif que la diffusion ne justifie pas. Sur un marché sans arbitrage, ce portage rémunère
ce que le modèle gaussien ne contient pas : la perte qu'un événement de crédit ou de liquidité infligerait au
porteur, ou des coûts de financement que le taux repo ne montre pas (décotes, coût de bilan). $F$ est alors
l'espérance du prix y compris cette perte, $M$ son espérance sur une trajectoire sans événement, et $\delta$ fait
passer de l'une à l'autre. L'approximation porte sur la loi : la moyenne du prix intègre la perte anticipée, sa
dispersion non, puisque la loi de $\bar B(T_0)$ reste celle d'un modèle gaussien, sans saut. Son effet reste limité tant que
$\delta$ est petit, ce que la section~\ref{sec:resultats} vérifie pour la courbe repo retenue.

\subsection{Prix exact}

Sous $\Q^{T_0}$, $\big(x(T_0), y(T_0)\big)$ est gaussien, de moyennes $(\mu_x, \mu_y + \delta)$, d'écarts-types
$(s_x, s_y)$ et de corrélation $r_{xy}$ (annexe~A). Le prix de l'OAT à l'échéance est une somme d'exponentielles
affines, $\bar B(T_0) = \sum_i K_iA_i\,e^{-H_x^i x(T_0) - H_y^i y(T_0)}$, décroissante en chaque facteur.
\citet[théorème~4.2.3]{BrigoMercurio2006} en tirent le prix exact d'une swaption, c'est-à-dire d'une option sur
obligation à coupons, sous forme d'intégrale à une dimension. On l'adapte en conditionnant par le facteur spread.
Sachant $y(T_0) = y$, $x(T_0)$ est gaussien de moyenne $m(y) = \mu_x + r_{xy}\,s_x\,(y - \mu_y - \delta)/s_y$ et
d'écart-type $s = s_x\sqrt{1 - r_{xy}^2}$, et le call est exercé si $x(T_0) < x^\ast(y)$, où $x^\ast(y)$ est
l'unique solution de $\bar B(T_0) = X$. Il vient
\begin{equation}
\begin{split}
C &= P(0,T_0)\int_{\mathbb R}\Big[\sum_i K_iA_i\,e^{-H_y^i y - H_x^i m(y) + \frac12 (H_x^i s)^2}\,\Phi\big(d(y) + H_x^i s\big)
- X\,\Phi\big(d(y)\big)\Big]\,\varphi_y(y)\,\dd y,\\
d(y) &= \frac{x^\ast(y) - m(y)}{s},
\end{split}
\label{eq:prix_exact}
\end{equation}
où $\varphi_y$ est la densité de $y(T_0)$ et $\Phi$ la fonction de répartition de la loi normale centrée réduite ;
le put s'obtient en remplaçant le crochet par $X\,\Phi(-d) - \sum_i(\cdots)\,\Phi(-d - H_x^i s)$. L'intégrale est
calculée par quadrature de Gauss-Hermite, et $x^\ast$ par la méthode de Newton sur $\log\bar B$. À $y$ fixé, on
retrouve la décomposition de Jamshidian : les zéro-coupons OAT sont des fonctions décroissantes du seul facteur
$x$, hypothèse sous laquelle \citet[CM3, p.~40-43]{Hillairet2026} l'établit dans un modèle HJM gaussien.

Le prix~\eqref{eq:prix_exact} adapte le théorème du livre sur trois points : il fait intervenir deux courbes
(l'actualisation €STR ne dépend que de $x$, l'OAT de $x + y$), les moyennes forward qui en découlent et le
recentrage $\delta$. Le livre intègre sur le premier facteur ; on intègre sur le second, ce qui garde l'intégrande
régulier quand $\sigma_y$ tend vers zéro et couvre ainsi le cas limite à un facteur, où le prix est celui de
Jamshidian.

\subsection{Approximation à poids gelés}

\citet{Russo} remplacent cette intégrale par une formule fermée, qui étend au G2++ l'approximation proposée à
un facteur par \citet{RussoFabozzi2016}. Si l'on gèle les poids $w_i$ des flux,
$\log\bar B(T_0) \approx c - \bar D_x\,x(T_0) - \bar D_y\,y(T_0)$ est gaussien, de variance
\begin{equation}
\bar\Sigma_B^2 = \bar D_x^2\,\Var x(T_0) + \bar D_y^2\,\Var y(T_0) + 2\bar D_x\bar D_y\,\Cov\big(x(T_0),y(T_0)\big),
\qquad \bar D_\cdot = \sum_{T_i>T_0} w_i\,H_\cdot(T_0,T_i),
\label{eq:sigma_b}
\end{equation}
et de moyenne fixée par $\E^{T_0}[\bar B(T_0)] = F$. \citet{Russo} écrivent cette variance comme une intégrale
évaluée numériquement ; avec des poids gelés, elle est fermée. Ils gèlent les poids à leur valeur d'aujourd'hui ;
on les gèle à leur valeur forward recentrée, $w_i = m_i(\delta)/F$, poids moyen de chaque flux dans le prix de
l'OAT à l'échéance. Le prix approché est alors celui de Black sur le forward repo :
\begin{equation}
C \approx P(0,T_0)\big[F\,\Phi(d_1) - X\,\Phi(d_2)\big],\qquad
d_{1,2} = \frac{\log(F/X) \pm \tfrac12\bar\Sigma_B^2}{\bar\Sigma_B}.
\label{eq:black}
\end{equation}
Cette formule serait exacte si le prix forward de l'OAT était une martingale lognormale sous $\Q^{T_0}$
\citep[CM3, p.~28]{Hillairet2026}. C'est le cas de chaque zéro-coupon, pas de leur somme : l'approximation est
bonne au centre de la distribution et se dégrade dans ses queues (section~\ref{sec:resultats}). Elle garde une
utilité même si l'on dispose du prix exact, car la décomposition de $\bar\Sigma_B^2$ en une part taux, une part
spread et une part croisée, de signe celui de $\rho$, sert à lire les résultats.
